\chapter{Posicionamiento Big Data y gobierno supervisado}\label{ap:posicionamiento}

Este anexo recoge dos discusiones complementarias que, por espacio, se han trasladado fuera del cuerpo de la memoria: el posicionamiento del trabajo frente a los entornos Big Data y la responsabilidad humana en el gobierno supervisado.

\section{Posicionamiento frente a entornos Big Data}\label{sec:posicionamiento-bigdata}

La descripción oficial propuesta para el TFE sitúa el trabajo en ``entornos Big Data y de computación en la nube'', expresión que conviene matizar para evitar una lectura equívoca. La plataforma no opera en la capa de procesamiento masivo de datos, sino en la capa de gobierno de metadatos e interoperabilidad. El caso de uso de validación trabaja con un volumen reducido de datos sintéticos y, por tanto, no constituye Big Data en el sentido estricto de volumen, velocidad y variedad. Esta frontera es coherente con la distinción establecida en el capítulo~\ref{cap:objetivos}: el objeto gobernado son los metadatos, no los datos de negocio.

La relevancia para entornos Big Data es, en consecuencia, indirecta pero sustantiva. Los ecosistemas de datos a gran escala se caracterizan por activos heterogéneos, distribuidos y de múltiples fuentes, cuya utilidad depende de que permanezcan descubribles, descritos y reutilizables. Esa función (catalogación, descripción semántica e interoperabilidad) es precisamente la que DAMA encuadra como gestión de metadatos dentro del gobierno del dato \cite{damaDmbok}. La plataforma demuestra dicho mecanismo de forma reproducible: descubre activos técnicos, los describe, los exporta como catálogo DCAT-AP-ES y valida su conformidad.

En un escenario Big Data real, el proyecto se utilizaría sin cambios en su lógica de negocio, dado que el flujo opera sobre metadatos y es independiente del volumen del dato subyacente. El crecimiento se absorbería por tres vías: (i) la incorporación de nuevas fuentes y servicios en OpenMetadata, catálogo concebido para inventarios técnicos amplios y heterogéneos \cite{openmetadataDocs}; (ii) la ampliación del contrato funcional añadiendo más tablas \texttt{gold} a la misma hoja de gobierno y al mismo workflow idempotente, sin modificar el núcleo; y (iii) el despliegue nativo en la nube, que habilita el escalado de la capa de catálogo y gobierno conforme a las primitivas de orquestación de Kubernetes \cite{kubernetesConcepts}. La salida interoperable en DCAT-AP-ES es, además, el vehículo por el que un patrimonio de datos extenso se hace federable y localizable entre portales nacionales y europeos \cite{datosGobDcatApEs2025,w3cDcat3}. En síntesis, el trabajo no escala procesando más datos, sino gobernando más metadatos, que es la dimensión en la que el gobierno aporta valor a un entorno Big Data.

\section{Gobierno supervisado y responsabilidad del operador}\label{sec:gobierno-supervisado}

Aunque el flujo es ampliamente automatizable y repetible, una parte esencial del gobierno reside en decisiones que deben permanecer bajo autoridad humana. Determinar qué activos son publicables, asignar la temática NTI-RISP \cite{datosGobNtiRisp}, clasificar un dataset como de alto valor o redactar su título, descripción y publicador no son transformaciones deterministas, sino juicios editoriales y de responsabilidad. El carácter manual y dinámico de esta curación no es una carencia del trabajo, sino una propiedad deseada del gobierno del dato.

El marco DAMA es explícito en este punto: el gobierno del dato asigna responsabilidad (\emph{accountability}) a roles definidos como el propietario y el \emph{steward} del dato, y la automatización debe asistir esa función, no sustituirla \cite{damaDmbok,datosGobDama2021}. Por consiguiente, la persona que opera o gobierna la plataforma debe conocer y aplicar las buenas prácticas del dominio, y responder de aquello que publica. La interoperabilidad técnica no exime de la responsabilidad sobre el contenido publicado.

Este principio tiene una implicación normativa concreta. La condición de conjunto de datos de alto valor está determinada por el Reglamento de Ejecución (UE) 2023/138 \cite{euHvdRegulation}; la plataforma trata la clasificación HVD como hipótesis de validación y no como una calificación jurídica automática, como ya se reconoció al discutir las limitaciones de los resultados. Por ello, la categorización debe ser revisada y autorizada por una persona responsable antes de cualquier publicación efectiva, y nunca derivarse de forma ciega por mera configuración.

El diseño de la plataforma materializa esta filosofía mediante salvaguardas explícitas. La asistencia de autorrellenado solo completa campos vacíos con valores de validación y nunca sobrescribe la curación introducida por el operador, de modo que la automatización ayuda sin desplazar el criterio humano. Del mismo modo, las líneas de trabajo futuro orientadas a la ejecución periódica mediante un planificador o CI/CD deben entenderse como automatización de la reproducibilidad, no como delegación de la decisión de publicar. La recomendación operativa es inequívoca: que el sistema sea técnicamente capaz de automatizar el flujo de extremo a extremo no implica que deba ejecutarse sin supervisión; la publicación de metadatos gobernados, por su visibilidad externa y por sus consecuencias jurídicas y reputacionales, debe ser un acto autorizado y supervisado.
